\subsection{Controlled Paired-Replay Extension}
\label{sec:paired_replay_results}

This experiment asks whether a matched healthy replay can provide useful detection-gated module evidence, and whether controlling sampling support improves the distinction between healthy and mutated runs. It evaluates the residual extension in Section~\ref{sec:paired_replay_method}, separately from the learned Stage 1--Stage 2--SHAP chain. In the corrected development comparison, the original learned chain detected 10/12 fault runs but also alarmed on 10/12 independent normal targets; no pair met the clean fault-specific detection criterion. Replacing its gate and then its module evidence with paired residuals yielded 10/12 detections, 1/12 normal alarms, Top-3 coverage of 10/12, and MRR 0.5833. The three-method comparison is provided in Supplementary Material Section S3. These development results motivated a frozen-parameter transfer evaluation rather than a reliability claim.

The main transfer result identifies a useful preprocessing improvement. On the same three transfer source flights, common-support residual computation reduced normal alarms from 11/12 to 3/12 under the frozen paired-residual gate, while retaining 10/12 fault detections. Top-3 and Top-5 both became 10/12, and MRR was 0.6667 (Table~\ref{tab:paired_transfer}). The common-support modification was examined after these transfer flights had been inspected and is therefore exploratory. It does not constitute a new independent confirmation, and its benefits cannot be attributed to retraining the original detector or to SHAP.

\begin{table}[htbp]
\centering\footnotesize
\caption{Paired residuals on the same transfer targets: frozen parameters and exploratory common-support processing.}
\label{tab:paired_transfer}
\renewcommand{\arraystretch}{1.15}
\setlength{\tabcolsep}{4pt}
\begin{tabular}{p{6.0cm}cc}
\toprule
Measure & \shortstack{Frozen paired\\residual frontend} & \shortstack{Common-support\\frontend}\\
\midrule
Fault detections & 12/12 & 10/12\\
Normal runs with an alarm & 11/12 & 3/12\\
Fault runs with pre-onset alarms & 10/12 & 3/12\\
Clean fault-specific detections & 1/12 & 7/12\\
Top-1 & 6/12 & 6/12\\
Top-3 & 9/12 & 10/12\\
Top-5 & 9/12 & 10/12\\
MRR & 0.6369 & 0.6667\\
Median detection delay (s) & 2.023 & 11.623\\
Maximum detection delay (s) & 54.819 & 107.623\\
\bottomrule
\end{tabular}
\par\smallskip\parbox{0.98\linewidth}{\scriptsize The threshold, feature noise and publisher allocation are unchanged. Both columns use 12 fault and 12 normal targets from the same three sources; normal alarms cover the complete available run. Delays are conditional on detection (12 and 10 detections, respectively); misses are retained in ranking denominators. The right column was evaluated after transfer data exposure.}
\end{table}


The two methods expose different detection trade-offs. The frozen reference method detected all 12 fault runs, but ten had already alarmed before injection and eleven corresponding normal targets alarmed. Common-support processing reduced the number of pre-injection fault alarms to three and increased clean fault-specific detections from one to seven. Its valid channel--run window fraction averaged 83.3\% (median 97.6\%); unavailable channel windows contributed no evidence. Thus, the comparison includes a change in usable observation support as well as the residual calculation, and coverage accompanies the reported improvement.

The remaining failures are concentrated by mutation. Under common support, Commander and Land Detector each produced 3/3 detections without paired normal alarms. All three EKF2 normal targets alarmed, and all three EKF2 fault runs alarmed before injection, so their 3/3 post-onset recall does not establish fault-specific detection. INAV produced only 1/3 detections, with a delay of 107.623 s; its other two trials remain misses. INAV and Land Detector rankings include tied publishers and do not identify a unique root cause. One Commander transfer trial failed the pre-existing effect-exclusivity check and remains in the 12-trial denominator; its score is not treated as an unambiguous diagnostic success. Supplementary Material Section S4 gives the per-mutation and per-run results.

Channel-specific calibration recovered stronger INAV candidate evidence, but the original-cache variant still produced 8/12 normal alarms despite Top-3/5 of 12/12. Further common-support calibration, within-flight calibration, and a two-healthy-reference variant did not resolve the overall recall--false-alarm trade-off. Their complete comparisons, including 24 additional healthy replays and a control on the same new normal targets, are retained in Supplementary Material Section S5. They are not combined with the best metrics of Table~\ref{tab:paired_transfer} to define a synthetic best method.
